\begin{table}[htbp]
\centering
\caption{Market variable calculation windows.}
\label{tab:mktvar_calc}
\small

\textbf{Panel A: MRET} --- trailing 104-week compounded market return (1,410 valid weeks)\\[4pt]
\begin{tabular}{r l l l}
\toprule
    Week & Window Start & Window End & MRET Date \\
\midrule
    104 & 1992-01-03 & 1993-12-24 & 1993-12-24 \\
    105 & 1992-01-10 & 1993-12-31 & 1993-12-31 \\
    106 & 1992-01-17 & 1994-01-07 & 1994-01-07 \\
    107 & 1992-01-24 & 1994-01-14 & 1994-01-14 \\
    108 & 1992-01-31 & 1994-01-21 & 1994-01-21 \\
    \multicolumn{4}{c}{$\vdots$} \\
    1511 & 2018-12-21 & 2020-12-11 & 2020-12-11 \\
    1512 & 2018-12-28 & 2020-12-18 & 2020-12-18 \\
    1513 & 2019-01-04 & 2020-12-25 & 2020-12-25 \\
\bottomrule
\end{tabular}
\vspace{12pt}

\textbf{Panel B: MVOL} --- trailing 24-month rolling standard deviation of monthly market returns (1,413 valid weeks)\\[4pt]
\begin{tabular}{r l}
\toprule
    Week & MVOL Date \\
\midrule
    101 & 1993-12-03 \\
    102 & 1993-12-10 \\
    103 & 1993-12-17 \\
    104 & 1993-12-24 \\
    105 & 1993-12-31 \\
    \multicolumn{2}{c}{$\vdots$} \\
    1511 & 2020-12-11 \\
    1512 & 2020-12-18 \\
    1513 & 2020-12-25 \\
\bottomrule
\end{tabular}
\vspace{4pt}
\begin{flushleft}
\footnotesize
\textit{Note.} MRET at week $t$ is the compounded gross market return (Mkt-RF + RF) over the preceding 104 weeks. MVOL at week $t$ is computed by first compounding weekly market returns to monthly frequency, then taking a 24-month rolling standard deviation. Each week inherits the value of its calendar month.
\end{flushleft}
\end{table}
